\documentclass[12pt]{article}
\usepackage[spanish]{babel}
\usepackage{amsmath, amssymb}
\usepackage[margin=2.5cm]{geometry}

\begin{document}
	
\section*{\centering Soluciones Guía 1}
	
	\section*{Parte I: Conjuntos numéricos}
	
	\begin{enumerate}
		\item $17+23+13+7 = 60 \in \mathbb{N}$
		\item $(12+18)+(8+2) = 40 \in \mathbb{N}$
		\item $6(15+5) = 120 \in \mathbb{N}$
		\item $8(20+10)-3(20+10) = 150 \in \mathbb{N}$
		\item $5(14+6)+4(14+6) = 180 \in \mathbb{N}$
		\item $7(12+8)+5(12+8)-3(12+8) = 180 \in \mathbb{N}$
	\end{enumerate}
	
	\section*{Parte II: Cierre de operaciones}
	
	\begin{enumerate}
		\item $8-13=-5 \notin \mathbb{N}$ $\Rightarrow$ No está cerrada en $\mathbb{N}$.
		\item $15\div5=3 \in \mathbb{Z}$, pero la división no es cerrada en $\mathbb{Z}$ en general.
		\item $\frac{2}{3}+\frac{5}{6}=\frac{3}{2} \in \mathbb{Q}$ $\Rightarrow$ Cerrada en $\mathbb{Q}$.
		\item $\sqrt{5}+3 \in \mathbb{R}$ $\Rightarrow$ Cerrada en $\mathbb{R}$.
		\item $7-12=-5 \in \mathbb{Z}$ $\Rightarrow$ Cerrada en $\mathbb{Z}$.
	\end{enumerate}
	
	\section*{Parte III: Propiedades}
	
	\begin{enumerate}
		\item $(17+23)+(13+7)=60$ (asociativa y conmutativa).
		\item $(12+18)+(8+2)=40$ (asociativa).
		\item $6(15+5)=120$ (distributiva).
		\item $(8-3)(20+10)=150$ (factor común).
		\item $(5+4)(14+6)=180$ (factor común).
		\item $(7+5-3)(12+8)=180$ (factor común).
	\end{enumerate}
	
	\section*{Parte IV: Análisis de error}
	
	\[
	(14-6)\cdot3 = 24
	\]
	\[
	14-(6\cdot3) = -4
	\]
	
	No son iguales, por lo tanto la afirmación es falsa.
	
	La propiedad correcta es:
	\[
	(a-b)c = ac - bc
	\]
	
	\section*{Parte V: Balance hídrico}
	
	\[
	18 + 11 - 23 - 4 = 2
	\]
	
	Balance total: $2$ mm.
	
	\[
	(18+11)-(23+4)=2
	\]
	
	Interpretación: ganancia de agua.
	
	\section*{Parte VI: Fertilización}
	
	\[
	150 \cdot 3.8 = 570 \text{ kg}
	\]
	
	\[
	570 \div 25 = 22.8 \Rightarrow 23 \text{ sacos}
	\]
	
	\section*{Parte VII: Mezclas}
	
	\[
	\text{Agua} = 2.4 \cdot 5 = 12 \text{ L}
	\]
	
	\[
	\text{Total} = 2.4 + 12 = 14.4 \text{ L}
	\]
	
	\section*{Parte VIII: Costo}
	
	\[
	4\cdot8200 + 3\cdot8200 + 2\cdot6900 = 71200
	\]
	
	\[
	(4+3)\cdot8200 + 2\cdot6900
	\]
	
	Costo total: \$71.200.
	
	\section*{Parte IX: Validación}
	
	\[
	140 \cdot 3.6 = 504 \text{ kg}
	\]
	
	\[
	504 \div 25 = 20.16 \Rightarrow 21 \text{ sacos}
	\]
	
\end{document}